% !TEX encoding = IsoLatin2  % notwendige Zeile f"ur Mac-Benutzer (muss als Kommentar stehen); Windos-Benutzer k"onnen die Zeile l"oschen.

% LaTeX-Vorlage Version 3.0,  Juni 2010
% erstellt von Dr. Andreas Drauschke und Dr. Susanne Teschl
% Modifiziert von Dr. Andreas Mehrle (andreas.mehrle@mci.edu)

\documentclass[a4paper,oneside,10pt]{article}%{scrreprt} 
% F"ur kurze Arbeiten w"are auch die Dokumentklasse "scrartcl" ausreichend. In diesem Fall ist "section" die h"ochste Ebene ("chapter" gibt es dann nicht).
% \documentclass[a4paper,bibtotoc,oneside]{scrbook}

\usepackage{ifthen}
\newboolean{english}
%\setboolean{english}{true}		% auskommentieren fuer deutsche Version

\ifthenelse{\boolean{english}}
{}{
	\usepackage{german}
}
\usepackage[dvips]{graphicx}
\usepackage[tmargin=1in,bmargin=1in,lmargin=1.25in,rmargin=1.25in]{geometry}
\usepackage{titlesec}
\usepackage{xcolor}
\usepackage[overload]{textcase}
\definecolor{MSBlue}{rgb}{.204,.353,.541}
\definecolor{MSLightBlue}{rgb}{.31,.506,.741}
% Set formats for each heading level
%\titleformat*{\section}{\rmfamily\bfseries\huge\color{MSBlue}\lowercase}
\titleformat{\section}[hang]{\rmfamily\bfseries\huge\color{MSBlue}\lowercase}{\thesection}{1em}{}[]
\titleformat{\subsection}{\large\bfseries\sffamily\uppercase}{\thesubsection}{1em}{}
\titleformat{\subsubsection}{\sffamily\bfseries}{\thesubsubsection}{1em}{}
\renewcommand\contentsname{\ifthenelse{\boolean{english}}{table of contents.}{inhaltsverzeichnis.}}
\renewcommand\listfigurename{\ifthenelse{\boolean{english}}{list of figures.}{abbildungsverzeichnis.}}
\renewcommand\listtablename{\ifthenelse{\boolean{english}}{list of tables.}{tabellenverzeichnis.}}
\renewcommand\refname{\litname.}
\usepackage{exscale,relsize}
\usepackage{fancyhdr}
\usepackage[small]{caption}
\usepackage{float}
\usepackage{amsmath}
\usepackage{units}
\usepackage{subfigure}
\usepackage{wallpaper}
\usepackage{pst-all}
\usepackage{amssymb}
\usepackage{rotating}
\usepackage[]{amsmath}
\usepackage{mathrsfs}
\usepackage{pgfplots}
\usepackage{makeidx}
%\usepackage[pdftex,bookmarks=true,bookmarksnumbered=true]{hyperref}

\newcommand{\babel}[2]{\ifthenelse{\boolean{english}}{#1}{#2}}

% color set
\definecolorseries{foo}{rgb}{last}[rgb]{1.0,0.0,0.0}[rgb]{0.0,0.0,1.0}
\resetcolorseries[16]{foo}

% auxillary symbols
\renewcommand{\tilde}{\symbol{126}}
\newcommand{\define}{\stackrel{!}{=}}
\renewcommand{\equiv}{\,\widehat{=}\,}
\newcommand{\subsubsubsection}{\textbf}
\newcommand{\re}{\mathrm{Re}}
\newcommand{\pr}{\mathrm{Pr}}
\newcommand{\st}{\mathrm{St}}
\newcommand{\fr}{\mathrm{Fr}}
\newcommand{\nus}{\mathrm{Nu}}
\newcommand{\gr}{\mathrm{Gr}}
\newcommand{\ra}{\mathrm{Ra}}
\newcommand{\mif}{\quad\mathrm{\babel{if}{falls}}\quad}
\newcommand{\with}{\quad\mathrm{\babel{with}{mit}}\quad}
\newcommand{\for}{\quad\mathrm{\babel{for}{f"ur}}\quad}
%\renewcommand{\not}{\not}
\newcommand{\im}{i}
\newcommand{\ariwam}{ARiWaM}
\newcommand{\matlab}{MATLAB}

% format specifications
\renewcommand{\emph}{\textbf}
\newcommand{\file}{\textit}
\newcommand{\cmd}{\texttt}
\newcommand{\ten}{\boldsymbol}
%\newcommand{\unit}{\mathrm}
\newcommand{\lemma}{\textit}
\newcommand{\deutsch}[1]{german: \textit{#1}}
\renewcommand{\index}{\emph}

% Command path to graphic files
\newcommand{\gpath}{./grafics}
\newcommand{\bsppath}{../uebungen/beispiele}

% mathematical operators
\newcommand{\grad}{\,\mathrm{grad}\,}
\renewcommand{\div}{\,\mathrm{div}\,}
\newcommand{\rot}{\,\mathrm{rot}\,}
\newcommand{\lap}{\Delta}
\newcommand{\laplace}[1]{\mathscr{L}\left\{#1\right\}}
\newcommand{\trans}{^T}
\newcommand{\norm}{\psarc[linewidth=0.5pt](0,0){0.4}{0}{90}\psdot[dotsize=0.1](0.15,0.15)}

% electrical networks
\newcommand{\con}{\pscircle[linewidth=2pt,fillstyle=solid,fillcolor=black](0,0){0.1}}
\newcommand{\pin}{\pscircle[linewidth=2pt,fillstyle=solid](0,0){0.15}}
\newcommand{\link}{\psline[linewidth=2pt]}
\newcommand{\linkd}{\psline[linestyle=dashed]}
\newcommand{\arcd}{\psarc[linestyle=dashed]}
\newcommand{\arrow}{\psline[linewidth=2pt,arrowsize=8pt]{->}}
\newcommand{\cs}{\pscircle[linewidth=2pt](0,0){0.5}\arrow(0,-0.4)(0,0.4)\link(0,-0.5)(0,-2)\link(0,0.5)(0,2)}
\newcommand{\vs}{\pscircle[linewidth=2pt](0,0){0.5}\psline(-0.2,-0.1)(0.2,-0.1)\psline(-0.2,0)(0.2,0)\pscurve(-0.2,0.1)(-0.1,0.15)(0,0.1)(0.1,0.05)(0.2,0.1)\link(0,-0.5)(0,-2)\link(0,0.5)(0,2)}
\newcommand{\res}{\psframe[linewidth=2pt](-0.3,-1)(0.3,1)\link(0,-1)(0,-2)\link(0,1)(0,2)}
\newcommand{\ind}{\psframe[linewidth=2pt,fillstyle=solid,fillcolor=black](-0.3,-1)(0.3,1)\link(0,-1)(0,-2)\link(0,1)(0,2)}
\renewcommand{\cap}{\link(-0.5,-0.1)(0.5,-0.1)\link(-0.5,0.1)(0.5,0.1)\link(0,-0.1)(0,-2)\link(0,0.1)(0,2)}
\newcommand{\opv}{\pspolygon[linewidth=2pt](-1,-1.3)(-1,1.3)(1.3,0)\link(-1,-0.8)(-2,-0.8)\link(-1,0.8)(-2,0.8)\link(1.3,0)(2,0)\rput(-0.7,0.8){$-$}\rput(-0.7,-0.8){$+$}}

% mechanical networks
\newcommand{\mass}{\psframe[linewidth=2pt](-0.5,-0.5)(0.5,0.5)\link(0,0.5)(0,1)\link(0,-0.5)(0,-1)}
\newcommand{\spring}{\link(0,-1)(0,-0.8)(-0.3,-0.7)(0.3,-0.5)(-0.3,-0.3)(0.3,-0.1)(-0.3,0.1)(0.3,0.3)(-0.3,0.5)(0.3,0.7)(0,0.8)(0,1)}
\newcommand{\damp}{\link(-0.5,0.4)(-0.5,-0.3)(0.5,-0.3)(0.5,0.4)\link(-0.42,0.2)(0.42,0.2)\link(0,0.2)(0,1)\link(0,-0.3)(0,-1)}
\newcommand{\hatch}{\psframe[fillstyle=vlines,linestyle=none,hatchwidth=0.25pt]}
\newcommand{\clamp}{\hatch(-0.5,-0.3)(0.5,0)\link(-0.5,0)(0.5,0)}
\newcommand{\fixed}{\hatch(-0.5,-0.7)(0.5,-0.4)\link(-0.5,-0.4)(0.5,-0.4)\link(-0.25,-0.4)(0,0)(0.25,-0.4)\pscircle[linewidth=2pt,fillstyle=solid](0,0){0.1}}
\newcommand{\loose}{\hatch(-0.5,-0.7)(0.5,-0.5)\link(-0.5,-0.4)(0.5,-0.4)\link(-0.5,-0.5)(0.5,-0.5)\link(-0.25,-0.4)(0,0)(0.25,-0.4)\pscircle[linewidth=2pt,fillstyle=solid](0,0){0.1}}
\newcommand{\axial}{\link(-0.3,-0.3)(-0.3,-0.1)(0.3,-0.1)(0.3,-0.3)\link(-0.3,0.3)(-0.3,0.1)(0.3,0.1)(0.3,0.3)}
\newcommand{\force}{\psline[linewidth=2pt,arrowsize=6pt]{->}}
\newcommand{\rmom}{\psarc[arrowsize=6pt,linewidth=2pt]{<-}}
\newcommand{\lmom}{\psarc[arrowsize=6pt,linewidth=2pt]{->}}
\newcommand{\stress}{\psline[arrowsize=6pt]{->}}
\newcommand{\beam}{\psframe[fillstyle=solid,fillcolor=lightgray]}
\newcommand{\rod}{\psline[linewidth=3pt]}
\newcommand{\disp}{\psline[linewidth=2pt,arrowsize=6pt]{->}}	% displacement

% hydraulical networks
\newcommand{\pump}{\pscircle[linewidth=2pt,fillstyle=solid,fillcolor=white](0,0){1}\pspolygon[fillstyle=solid,fillcolor=black](0,1)(-0.15,0.3)(0.15,0.3)}
\newcommand{\motor}{\pscircle[linewidth=2pt,fillstyle=solid,fillcolor=white](0,0){0.8}\rput(0,0){M}}
\newcommand{\throttle}{\psarc[linewidth=2pt](-3,0){2.7}{-15}{15}\psarc[linewidth=2pt](3,0){2.7}{165}{195}}
\newcommand{\limiter}{\psframe[linewidth=2pt,fillstyle=solid,fillcolor=white](-0.7,-0.7)(0.7,0.7)\psline{->}(-0.5,-0.4)(-0.5,0.1)(0.5,-0.25)}

% diagrams (block and s-plane)
\newcommand{\axis}{\psline[arrowsize=6pt,arrowinset=0]{->}}
\newcommand{\measure}{\psline[linewidth=0.5pt,arrowsize=4pt]}
\newcommand{\xtick}[1]{\measure(#1,-0.1)(#1,0.1)}
\newcommand{\ytick}[1]{\measure(-0.1,#1)(0.1,#1)}
\newcommand{\ang}{\psarc[linewidth=0.5pt,arrowsize=4pt]{<->}}
\newcommand{\tf}[1]{\psframe[linewidth=2pt](-1,-0.75)(1,0.75)\rput(0,0){#1}}
\newcommand{\tff}[2]{\psframe[fillstyle=solid,fillcolor=lightgray](-1,0)(1,0.75)\psframe[linewidth=2pt](-1,-0.75)(1,0.75)\rput(0,0.35){#1}\rput(0,-0.35){#2}}
\newcommand{\sat}{\psframe[linewidth=2pt](-1,-0.75)(1,0.75)\link(-0.8,-0.5)(-0.5,-0.5)(0.5,0.5)(0.8,0.5)}
\newcommand{\add}[4]{\pscircle[linewidth=2pt,fillstyle=solid,fillcolor=white](0,0){0.5}\rput(-0.5,-0.75){#1}\rput(-0.75,0.5){#2}\rput(0.5,0.75){#3}\rput(0.75,-0.5){#4}}
\newcommand{\pole}{\link(-0.15,-0.15)(0.15,0.15)\link(-0.15,0.15)(0.15,-0.15)}
\newcommand{\zero}{\pscircle[linewidth=2pt,fillstyle=solid](0,0){0.15}}
\newcommand{\textblock}[1]{\psframebox[linewidth=2pt]{\parbox[c][2cm][c]{3cm}{\centering #1}}}

% thermodynamic systems
\newcommand{\cv}{\psframe[linestyle=dashed]}
\newcommand{\sys}{\psframe[fillcolor=white,fillstyle=solid,linewidth=2pt]}
\newcommand{\iso}{\psframe[linewidth=2pt,fillstyle=vlines]}
\newcommand{\candle}{\psframe[linewidth=0.5pt](-0.15,-0.5)(0.15,0.3)\pscurve[linewidth=0.5pt](0,0.8)(-0.1,0.5)(0,0.3)(0.1,0.5)(0,0.8)\pscurve[linewidth=0.5pt](0,0.6)(-0.08,0.4)(0,0.3)(0.08,0.4)(0,0.6)}
\newcommand{\flux}{\psline[arrowsize=10pt,linewidth=4pt]{->}}
\newcommand{\point}[1]{\pscircle(0,0){0.3}\rput[c](0,0){#1}}

% aerodynamics
\newcommand{\source}{\pin\rput{0}(0,0){\arrow(0.5,0)(1,0)}\rput{45}(0,0){\arrow(0.5,0)(1,0)}\rput{90}(0,0){\arrow(0.5,0)(1,0)}\rput{135}(0,0){\arrow(0.5,0)(1,0)}\rput{180}(0,0){\arrow(0.5,0)(1,0)}\rput{225}(0,0){\arrow(0.5,0)(1,0)}\rput{270}(0,0){\arrow(0.5,0)(1,0)}\rput{315}(0,0){\arrow(0.5,0)(1,0)}}

% CFD
\newcommand{\cell}[1]{\definecolor{cellc}{gray}{#1}\psframe[linewidth=2pt,fillstyle=solid,fillcolor=cellc](0,0)(1,1)}

%\renewcommand{\labelenumi}{\alph{enumi})}

\setlength{\parindent}{0em}
\setlength{\parskip}{1.5ex plus0.5ex minus0.5ex}
\setlength{\captionmargin}{3em}

% counters
\newcounter{example}
\ifthenelse{\boolean{english}}{\newcommand{\exampletext}{example }}{\newcommand{\exampletext}{Beispiel }}
\newcommand{\example}[1]{\underline{\exampletext \arabic{chapter}.\arabic{example}:} #1 \addtocounter{example}{1}}
\newcounter{exercise}
\ifthenelse{\boolean{english}}{\newcommand{\exercisetext}{exercise }}{\newcommand{\exercisetext}{Aufgabe }}
\newcommand{\exercise}[1]{\underline{\exercisetext \arabic{chapter}.\arabic{exercise}:} #1 \stepcounter{exercise}}
\newcommand{\cchapter}[1]{\chapter{#1} \setcounter{example}{1} \setcounter{exercise}{1}}
\ifthenelse{\boolean{english}}{\newcommand{\solution}{\textit{solution:} }}{\newcommand{\solution}{\textit{L\"osung:} }}
%\newcounter{beispiel}
%\setcounter{beispiel}{1}		% Nummer des ersten Beispiels
\newboolean{student}
\newcounter{enumcount}
\newcommand{\resume}[1]{\begin{#1} \setcounter{enumi}{\value{enumcount}}}
\newcommand{\pause}[1]{\setcounter{enumcount}{\value{enumi}} \end{#1}}




\def\labtitle{Labor Biosignalverarbeitung}
\def\labcode{MTM-B-5-ATR-ASA-LB}
\def\labname{Signalverarbeitung von Biosignalen}
\def\labnum{1}

\def\study{Master Mechatronik}
\def\term{2.}
\def\lecturer{Andreas Mehrle}
\def\group{[Gruppe und Gruppennummer]}
\def\student{[Namen der Gruppenmitglieder]}

\def\litname{literatur}

\begin{document}

%\layout

% Kopf- und Fusszeilen initiieren
\pagestyle{plain}

% Deckblatt:
\thispagestyle{empty}
\sffamily
%\put(-30,-685){\includegraphics[width=1.15\linewidth]{BG}}

{\bfseries


\vspace*{1cm}

\def\thesis{\ifthenelse{\boolean{english}}{laboratory report}{laborbericht}}
\textcolor{MSBlue}{\rmfamily\bfseries\huge\thesis.}%

\vspace*{0.5cm}

\textbf{\labtitle{ }(\labcode)}

\labname

\ifthenelse{\boolean{english}}{Lab}{Labor} \labnum

\vspace*{10cm}

\textcolor{gray}{\study}

\def\termname{\ifthenelse{\boolean{english}}{semester}{Semester}}
\textcolor{gray}{\term{ }\termname}

\def\lecturername{\ifthenelse{\boolean{english}}{Responsible lecturer}{Lehrveranstaltungsleiter}}
\textcolor{gray}{\lecturername: \lecturer}

\def\groupname{\ifthenelse{\boolean{english}}{Group}{Jahrgang}}
\textcolor{gray}{\groupname: \group}

\def\authorname{\ifthenelse{\boolean{english}}{Author}{Verfasser}}
\textcolor{gray}{\authorname: \student}

\textcolor{gray}{\today}

\newpage}


\tableofcontents%\thispagestyle{empty}

\newpage

% Kopf- und Fusszeilen initiieren
\pagestyle{fancy}

\section{Aufgabe 1: Darstellung des Messsystems}
\label{sec:messsystem}

Abbildung~\ref{fig:messsystem} zeigt das verwendete Messsystem. Beschriften Sie alle Komponenten, Signalpfade und Busse einschliesslich Bustyp und Geschwindigkeit. Beschreiben Sie jede Komponente in ein bis zwei S"atzen und kennzeichnen Sie eigene oder "ubernommene Darstellungen entsprechend.

\begin{figure}[H]
\centering
\fbox{\parbox[c][4cm][c]{0.85\textwidth}{\centering
[Blockdiagramm des Messsystems hier einf"ugen]}}
\caption{Messsystem mit Komponenten, Signalpfaden und beschrifteten Bussen.}
\label{fig:messsystem}
\end{figure}

\section{Aufgabe 2: Umrechnung des Temperatursensors}
\label{sec:temperatur}

Der Sensor deckt $0$--$3{,}3\,\mathrm{V}$ beziehungsweise $0$--$40\,^\circ\mathrm{C}$ ab. Der 10-Bit-ADC verwendet einen Messbereich von $0$--$5\,\mathrm{V}$. Berechnen Sie aus dem ADC-Wert 520 zun"achst die Eingangsspannung und daraus die Temperatur. Dokumentieren Sie Rechenweg, verwendete ADC-Konvention und Ergebnis:
\begin{align}
U_\mathrm{ADC} &= \frac{\text{ADC-Wert}}{2^{10}-1}\,U_\mathrm{ref},\\
T &= \frac{U_\mathrm{ADC}}{3{,}3\,\mathrm{V}}\,40\,^\circ\mathrm{C}.
\end{align}
\noindent
[Rechnung und Ergebnis einf"ugen.]

\section{Aufgabe 3: Achsen und Einheiten des Beschleunigungssensors}
\label{sec:achsen}

Ordnen Sie die drei ausgegebenen Beschleunigungswerte den Sensorachsen zu und beschreiben Sie einen Test zur Zuordnung. Geben Sie die Einheit der Rohdaten an und erl"autern Sie anhand des Sensordatenblatts die Umrechnung in $\mathrm{m/s^2}$. Nennen Sie Datenblatt, Seite, Messbereich und verwendeten Umrechnungsfaktor.

[Achsen, Test, Einheit und Umrechnung anhand des verwendeten Sensortyps dokumentieren.]

\section{Aufgabe 4: Vertikale Handbewegung mit Standardeinstellung}
\label{sec:handbewegung}

\subsection{Messung und Plot (Aufgabe 4a)}
Zeigen Sie vier Sekunden der zehn Sekunden langen Messung. Achsen und Einheiten m"ussen beschriftet sein; verweisen Sie im Text auf die Abbildung.

\begin{figure}[H]
\centering
\fbox{\parbox[c][4cm][c]{0.85\textwidth}{\centering
[Plot der ersten Messung "uber vier Sekunden hier einf"ugen]}}
\caption{Beschleunigungsdaten bei schneller vertikaler Handbewegung mit Standardeinstellung.}
\label{fig:handbewegung-standard}
\end{figure}

\subsection{Plateaus der Peaks (Aufgabe 4b)}
Diskutieren Sie, weshalb die Peaks Plateaus mit demselben Maximal- oder Minimalwert aufweisen und wie sich dieses Verhalten vermeiden l"asst.

[Beobachtung und Erkl"arung einf"ugen.]

\subsection{Abweichung von der ausgef"uhrten Bewegung (Aufgabe 4c)}
Erl"autern Sie, weshalb das gemessene Signal die ausgef"uhrte Bewegung nicht exakt abbildet.

[Diskussion einf"ugen.]

\section{Aufgabe 5: Messung mit angepasster Sensorkonfiguration}
\label{sec:sensorkonfiguration}

Dokumentieren Sie die "Anderung am Arduino-Code und die zugeh"orige Information aus dem Sensordatenblatt: Seite, Register, Parameter und alle zul"assigen Parameterwerte. Begr"unden Sie die gew"ahlte Einstellung. Zeigen Sie anschliessend den Plot einer wiederholten Messung "uber vier Sekunden und vergleichen Sie ihn mit Abbildung~\ref{fig:handbewegung-standard}.

\begin{figure}[H]
\centering
\fbox{\parbox[c][4cm][c]{0.85\textwidth}{\centering
[Plot der Messung mit angepasster Sensorkonfiguration hier einf"ugen]}}
\caption{Beschleunigungsdaten nach Anpassung der Sensorkonfiguration.}
\label{fig:handbewegung-angepasst}
\end{figure}

[Gemeinsamkeiten und Unterschiede der beiden Plots sowie deren Ursachen diskutieren.]

\section{Aufgabe 6: Abtastfrequenz und Messung der sechs Sensorlagen}
\label{sec:abtastung}

Beschreiben Sie, wie jeder Messwert unmittelbar auf dem Mikrocontroller mit einem Zeitstempel versehen wurde. Geben Sie die Berechnung der mittleren Abtastfrequenz und ihrer Abweichung aus den Zeitstempeln an. Dokumentieren Sie die Messungen von jeweils zehn Sekunden f"ur die sechs Lagen (jede Achse nach oben und unten) und pr"ufen Sie, ob jeweils etwa $1\,g$ angezeigt wird.

\subsection{Plot der sechs Lagen (Aufgabe 6a)}
\begin{figure}[H]
\centering
\fbox{\parbox[c][4cm][c]{0.85\textwidth}{\centering
[Plot der Messungen in allen sechs Sensorlagen hier einf"ugen]}}
\caption{Beschleunigungsdaten in den sechs Sensorlagen.}
\label{fig:sechs-lagen}
\end{figure}

\subsection{Plausibilit"at der Achsenwerte (Aufgabe 6b)}
Diskutieren Sie, ob die gemessenen Werte bei nach oben beziehungsweise nach unten gerichteter Achse physikalisch plausibel sind. Ber"ucksichtigen Sie Vorzeichen, Lage und Erdbeschleunigung.

[Beobachtungen und Plausibilit"atspr"ufung einf"ugen.]

\section{Aufgabe 7: Wirkung des Filters}
\label{sec:filter}

Stellen Sie gefilterte und ungefilterte Daten aus dem vorherigen Experiment gemeinsam dar. Beschreiben Sie in ein bis zwei S"atzen die Filterwirkung und begr"unden Sie die Wahl des Filters.

\begin{figure}[H]
\centering
\fbox{\parbox[c][4cm][c]{0.85\textwidth}{\centering
[Vergleich von gefilterten und ungefilterten Daten hier einf"ugen]}}
\caption{Ungefilterte und gefilterte Beschleunigungsdaten.}
\label{fig:filtervergleich}
\end{figure}

[Kurze Beschreibung der Filterwirkung und Begr"undung einf"ugen.]

\section{Aufgabe 8: Bestimmung der Schrittfrequenz}
\label{sec:schrittfrequenz}

Bestimmen Sie aus den aufgezeichneten Daten die mittlere Schrittfrequenz und ihre Abweichung. Erl"autern Sie kurz die Erkennung der Schritte und die Berechnung.

[Methode, Mittelwert, Abweichung und Ergebnis einf"ugen.]

\section{Aufgabe 9: Vergleich mit dem mobilen Messsystem}
\label{sec:mobil}

Nehmen Sie mit Lab1Code2, DataLogger und Batterie etwa zehn Sekunden auf. Vergleichen Sie Mittelwert und Abweichung der Abtastfrequenz mit Abschnitt~\ref{sec:abtastung}. Erkl"aren Sie in einem Absatz die beobachteten Unterschiede.

[Messbedingungen, Vergleichswerte und Erkl"arung einf"ugen.]

\section{Aufgabe 10: Funktionspr"ufung des Python-Codes}
\label{sec:zusammenfassung}

Pr"ufen und dokumentieren Sie, dass \file{MainCode1.ipynb} mit den mitgelieferten Datens"atzen l"auft und die in den Aufgaben geforderten Plots erzeugt. Fassen Sie die wichtigsten Ergebnisse und Erkenntnisse anschliessend kurz zusammen.

\section*{Hinweis zur Codeabgabe}
Das ausf"uhrbare Hauptskript der Python-Auswertung ist separat als \file{MainCode1.ipynb} abzugeben. Legen Sie ausserdem den modifizierten Arduino-Code, alle zum Ausf"uhren ben"otigten Datens"atze und selbstgeschriebene Python-Bibliotheken in die geforderte ZIP-Datei. Der Bericht darf einschliesslich Abbildungen und Verzeichnissen h"ochstens zehn Seiten umfassen.

\bibliographystyle{IEEEtran}	
\bibliography{Literatur}

\end{document}
